\section{Capability Evaluation}

\subsection{Experimental Setup}

\paragraph{Models.}
We evaluate eight advanced LLMs on FinIF-Test, spanning two categories.
\textit{Closed-source models}: GPT-5.5 and GPT-5~\cite{openai2025gpt5}.
\textit{Open-weighted models}: GLM-5.1~\cite{glm2025chatglm}, DeepSeek-V4-Pro and DeepSeek-V4-Flash~\cite{deepseek2025v4}, Qwen3.5-27B, Qwen3.5-9B, and Qwen3.5-4B~\cite{qwen2025qwen3}.
For all models, we set the sampling temperature to 0 to ensure deterministic outputs.

\paragraph{Metrics.}
We adopt two primary evaluation metrics.
\textbf{Instruction-level Success Rate (ISR)} measures the fraction of test instructions whose \emph{every} annotated constraint is satisfied:
\begin{equation}
\text{ISR} = \frac{1}{N}\sum_{i=1}^{N} \mathbb{1}\!\left[\,\forall\, c \in \mathcal{C}_i,\ \text{pass}(c) = 1\right]
\label{eq:isr}
\end{equation}
where $N$ is the number of test instructions and $\mathcal{C}_i$ is the constraint set for instruction $i$.
ISR captures whether a model's output can serve as a complete, audit-ready work product: even a single violated constraint renders the entire instruction unsuccessful.
\textbf{Quality Score} is a holistic 0--10 score assigned by a calibrated LLM judge, assessing the financial validity, reasoning soundness, and professional completeness of each response as an auxiliary diagnostic.

\paragraph{Hybrid Verification Engine.}
FinIF-Test contains 3,282 fine-grained constraints, comprising 1,401 structural constraints evaluated by deterministic Python rule checkers and 1,881 semantic constraints judged by GPT-4o.
The evaluator follows a \emph{rule-first} protocol: mechanically checkable constraints (e.g., word counts, decimal-place precision, table column presence, required headings, ordered-list formats) are routed to deterministic checkers, while semantic constraints (e.g., evidence grounding, decision boundary reasoning, calculation lineage) are batched per item and judged by GPT-4o with fixed prompts, temperature 0, and strict JSON output parsing.
This hybrid design ensures reproducibility for structural checks while preserving coverage for constraints requiring semantic interpretation.

\subsection{Main Results}

\begin{table}[t]
\centering
\small
\resizebox{\linewidth}{!}{
\begin{tabular}{lrrrr}
\toprule
Model & Micro IF & Macro IF & Strict & Quality \\
\midrule
GPT-5.5 & 93.30 & 93.10 & 58.67 & 7.40 \\
GPT-5 & 91.35 & 91.20 & 47.33 & 7.05 \\
DS-V4-Pro & 81.24 & 80.72 & 16.00 & 4.68 \\
DS-V4-Flash & 77.92 & 77.25 & 9.00 & 4.27 \\
Qwen3.5-9B & 72.08 & 71.45 & 4.00 & 3.60 \\
\bottomrule
\end{tabular}
}
\caption{Main results on FinIF-Test. IF metrics are percentages. Quality is a 0--10 auxiliary judge score. Strict pass requires satisfying every IF constraint in an item.}
\label{tab:main_results}
\end{table}

Table~\ref{tab:main_results} shows a consistent gap between local constraint compliance and complete work-product compliance.
GPT-5.5 passes 2,924 of 3,134 IF constraints, but strictly passes only 176 of 300 items.
GPT-5 passes 2,863 constraints but strictly passes only 142 items.
The gap is larger for cheaper and smaller models: DS-V4-Pro strictly passes 48 items, DS-V4-Flash 27 items, and Qwen3.5-9B only 12 items.

This gap is the central empirical finding of FinIF.
Micro IF captures whether a model follows many local requirements, while strict item pass rate captures whether the response can be accepted as a complete financial work product.
In production-like workflows, a response that violates one critical evidence, calculation, or boundary instruction may still require human repair even if most local constraints pass.

\subsection{Further Analysis}

\paragraph{Evidence placement is the dominant failure mode.}
Across the evaluated models, the most frequent failed tag is EG2, which requires active source labels next to material facts, figures, thresholds, missing-evidence triggers, and decision points.
GPT-5 fails 119 EG2 constraints, and GPT-5.5 fails 107.
This indicates that even strong models often cite sources globally but fail to place evidence exactly where an audit trail requires it.

\paragraph{Quantitative verification remains fragile.}
The second recurring weakness is QV2, which requires source inputs, formula or comparison, result, and business implication for calculations and reconciliations.
GPT-5 fails 42 QV2 constraints, while DS-V4-Pro and Qwen3.5-9B fail 113 and 162, respectively.
These failures are not necessarily arithmetic mistakes; many are lineage failures where the response gives a number without the required verification trail.

\paragraph{Decision-boundary reasoning is hard to keep connected.}
DB9 asks the response to connect active evidence, governing rule or task standard, and final action in a single reasoning chain.
GPT-5.5 and GPT-5 still fail 27 and 22 DB9 constraints, respectively.
Smaller models fail this tag much more often.
This suggests that models can list facts and recommendations, but often struggle to maintain an auditable evidence-rule-action chain.

\paragraph{Format constraints are not solved.}
Although many format checks are mechanically simple, they continue to cause item failures.
GPT-5.5 fails 19 FP3 constraints, and GPT-5 fails 58.
These include length, decimal-place, count, and ordering requirements.
The Figure~\ref{fig:case} example shows why this matters: an otherwise high-quality response can be rejected because it fails a concrete delivery constraint.

\section{FinIF-Train for Model Adaptation}
\label{sec:sft_experiment}

FinIF-Train provides the training component of FinIF.
It contains 1,064 workflow-grounded examples across 38 task types and 7,851 annotated constraints.
Unlike generic instruction-tuning data, each example is tied to a finance workflow, a concrete work product, and explicit delivery constraints.
The intended use is supervised fine-tuning or preference-data construction for financial instruction following.

To avoid conflating dataset versions, we do not report legacy pilot SFT numbers in this draft.
The same-split adaptation protocol is straightforward: train on FinIF-Train, evaluate both the base model and the adapted model on the 300-item FinIF-Test split, and report Micro IF, Macro item IF, Strict item pass rate, and auxiliary quality.
This keeps model adaptation results directly comparable with the capability evaluation above.
